% Appendix: the logit-distance test. Negative, and in the appendix because it
% is a statement about a proxy rather than about the theory it tests.
\section{Does Logit Distance Predict Probe Transportability?}
\label{sec:appendix:logit}

\citet{nielsen2026logit} identify a logit-based distributional distance as the
quantity that controls linear representational similarity, and show that KL
closeness does not. Two models can match each other's predictions while failing
to preserve the linear recoverability of the same concepts. That makes logit
distance the better-motivated predictor of whether a probe survives the map
between two models, where Section~\ref{sec:results} reports CKA.

We can test a version of it at no cost, because extraction caches each model's
out-loud probability that a claim is true. For models $A$ and $B$ on the same
items,
\[
  d_{\mathrm{logit}} = \frac{1}{n}\sum_i
  \left| \log\frac{p^A_i}{1-p^A_i} - \log\frac{p^B_i}{1-p^B_i} \right| ,
\]
alongside the mean KL between the two Bernoulli distributions. We correlate each
against the transfer loss of Section~\ref{sec:results}, over all $208$ cells.

\begin{table}[h]
\centering
\small
\begin{tabular}{@{}lrrr@{}}
\toprule
 & \multicolumn{3}{c}{correlation with transfer loss} \\
\cmidrule(lr){2-4}
predictor & $A\!\to\!B$ & $B\!\to\!A$ & worse direction \\
\midrule
logit distance (mean absolute) & $-0.010$ & $0.036$ & $0.051$ \\
KL                             & $-0.057$ & $0.083$ & $0.050$ \\
correlation of $p(\text{yes})$ & $-0.451$ & $-0.278$ & $-0.487$ \\
CKA                            & $\mathbf{-0.574}$ & $-0.151$ & $\mathbf{-0.514}$ \\
map $R^2$ ($A\!\to\!B$)        & $-0.388$ & $0.070$ & $-0.301$ \\
\bottomrule
\end{tabular}
\caption{Negative correlations mean the predictor anticipates \emph{better}
transfer. Pooled over $208$ cells.}
\label{tab:logit-pooled}
\end{table}

\begin{table}[h]
\centering
\small
\begin{tabular}{@{}lrrrr@{}}
\toprule
dataset & logit distance & KL & CKA & $n$ \\
\midrule
geometry\_of\_truth & $+0.306$ & $+0.515$ & $-0.598$ & 45 \\
boolq               & $+0.301$ & $+0.110$ & $-0.522$ & 45 \\
truthfulqa          & $+0.185$ & $+0.310$ & $-0.525$ & 45 \\
imdb                & $+0.041$ & $+0.350$ & $-0.663$ & 28 \\
rte                 & $-0.233$ & $-0.018$ & $-0.139$ & 45 \\
\bottomrule
\end{tabular}
\caption{Within each dataset, against the worse of the two transfer directions.
For logit distance a \emph{positive} value is the predicted direction, so more
distance, more loss.}
\label{tab:logit-bydataset}
\end{table}

Pooled, the proxy predicts nothing, and CKA remains the better available
predictor. Within each dataset the sign is the predicted one in four cases of
five, so the direction is present and the magnitude is weak; the pooled estimate
vanishes because variation between datasets swamps it.

We read this as a statement about the proxy and not about the result it tests.
The distance in \citet{nielsen2026logit} is defined over a model's full
conditional distribution, and ours projects that onto a single binary question --
one dimension of a vocabulary-sized object -- which is evidently too narrow to
carry the quantity. Testing the theory in this setting requires caching full
next-token distributions over a corpus, which we did not do; that is the concrete
addition for anyone who wants the answer. Consistently with their account, KL is
equally uninformative here.
